\section{Datasets} \label{sec:datasets}

This chapter describes the datasets and the data selection filters applied in the validation runs of this report.
\\[1em]
\noindent The validation period is from \textbf{\detokenize{$<ReportVars['Common']['interval_from']>$} to \detokenize{$<ReportVars['Common']['interval_to']>$}}.
Only measurements within this period are used for validation.

\subsection{Candidates -- \ac{C3S} SM COMBINED, ACTIVE and PASSIVE (\detokenize{$<next(v for v in Run1ContentVars['ConfigVars'].values() if isinstance(v, dict) and v.get('name') == 'C3S SM combined')['version']>$})} \label{ssec:c3s}

\noindent The \ac{C3S} satellite soil moisture products (\url{https://doi.org/10.24381/cds.d7782f18}) are global, long-term \acp{CDR} of daily surface soil moisture (representative of the top few centimetres of the soil) on a regular 0.25° grid \citep{c3satbd}.
The algorithm is based on that of the \ac{ESA} \ac{CCI} Soil Moisture project \citep{liu2012,dorigo2017,gruber2019}.
Three products are provided:
\begin{itemize}
    \item \textbf{ACTIVE}: merges soil moisture retrievals from active microwave sensors (scatterometers, e.g.\ \ac{ASCAT}). It is provided as degree of saturation (in \%).
    \item \textbf{PASSIVE}: merges soil moisture retrievals from passive microwave sensors (radiometers, e.g.\ \ac{SMOS}, \ac{SMAP}, AMSR2, GPM), in \unitsm{}.
    \item \textbf{COMBINED}: merges the active and passive retrievals into a single harmonised record, in \unitsm{}, available from November 1978 to present.
\end{itemize}
For the COMBINED product, all sensor time series are first rescaled into a common climatology (via CDF matching against \ac{GLDAS} Noah land surface model simulations), their random errors are estimated via \ac{TCA}, and the rescaled observations are then merged using weighted averaging, with weights based on the estimated errors.
A final break detection and correction step reduces inhomogeneities introduced by changes in the sensor constellation over time \citep{preimesberger2021}.
\\[1em]
\noindent The \acp{CDR} are updated once per year and extended in between by an \ac{ICDR}, which is based on the same input data streams and algorithm and therefore has very similar quality characteristics \citep{c3spqar}.
\ac{ICDR} extensions are published every 10 days, each containing 10 days of new data, with a delay of 1--2 weeks after the last day covered by the extension.
Data are available through the Copernicus Climate Data Store; the full quality assessment of each product version is published in the \ac{PQAR} \citep{c3spqar}.
In \ac{QA4SM}, the records are regularly extended with new \ac{ICDR} data.
\\[1em]
\noindent The \ac{C3S} COMBINED product targets an accuracy (\ac{ubRMSD}) of 0.01--\kpithres{}~\unitsm{}, depending on land cover and climate (\ac{KPI} KPI.SM.D3.1, \citealp{c3spqar}); \ac{GCOS} specifies \gcosthres{}~\unitsm{} for surface soil moisture \citep{gcos2022}.
Known limitations apply over dense vegetation, frozen or snow-covered soils (where no soil moisture is provided), and in regions affected by radio-frequency interference \citep{dorigo2017}.

\subsubsection{Data filters}

\noindent All filters are applied \textit{before} quality metrics are computed.
For a detailed description of each filter, see the validation run online (link in Table \ref{tab:val_run_list}) and the \href{https://qa4sm.eu/api/user-manual}{QA4SM User Manual}.
The following \textit{binary} filters were activated (the same in both validation runs):
\begin{itemize}
    \item COMBINED: \detokenize{$<', '.join([d['description'] for d in next(v for v in Run1ContentVars['ConfigVars'].values() if isinstance(v, dict) and v.get('name') == 'C3S SM combined')['basic_filters_description']])>$}
    \item ACTIVE: \detokenize{$<', '.join([d['description'] for d in next(v for v in Run1ContentVars['ConfigVars'].values() if isinstance(v, dict) and v.get('name') == 'C3S SM active')['basic_filters_description']])>$}
    \item PASSIVE: \detokenize{$<', '.join([d['description'] for d in next(v for v in Run1ContentVars['ConfigVars'].values() if isinstance(v, dict) and v.get('name') == 'C3S SM passive')['basic_filters_description']])>$}
\end{itemize}

% ---------------------------------------------------------------------------------------------------------------------

\subsection{Reference -- ISMN in situ measurements (\detokenize{$<next(v for v in Run1ContentVars['ConfigVars'].values() if isinstance(v, dict) and v.get('name') == 'ISMN')['version']>$})} \label{ssec:ismn}

\noindent The \ac{ISMN} \citep{dorigo2021} is a global, public collection of in situ soil moisture measurements from operational and historic networks and measurement campaigns\footnote{available for download after registration at \url{https://ismn.earth}}.
The \ac{ISMN} collects, harmonizes and publishes ground measurements of soil moisture and related geophysical variables from various network maintainers across the globe.
Data within the \ac{ISMN} are subject to quality controls and provided with quality flags \citep{dorigo2013}.
Some networks in Europe and the United States provide data in \ac{NRT}, meaning that new measurements are automatically added to the \ac{ISMN} database with a short delay.
\\[1em]
\ac{QA4SM} hosts multiple downloaded snapshots of the full \ac{ISMN} archive.
The version used in this report is \detokenize{$<next(v for v in Run1ContentVars['ConfigVars'].values() if isinstance(v, dict) and v.get('name') == 'ISMN')['version']>$}.
\ac{NRT} networks in this version are regularly updated in \ac{QA4SM} and are the only networks used in this report, i.e., non-NRT networks are currently not used.
In line with the \ac{PQAR} \citep{c3spqar}, all sensors in 0--10\,cm depth are used.
In addition, only sensors classified as \ac{FRMs} are used, i.e.\ the subset of \ac{ISMN} sensors that is recommended for satellite validation; see the \ac{PQAR} for their definition \citep{c3spqar}.

\subsubsection{Data filters}

\noindent The following \textit{binary} filters were activated for \ac{ISMN}:
\begin{itemize}
\detokenize{$<' '.join([f"\\item {d['description']}" for d in next(v for v in Run1ContentVars['ConfigVars'].values() if isinstance(v, dict) and v.get('name') == 'ISMN')['basic_filters_description']])>$}
\end{itemize}
\vspace{1em}
\noindent The following \textit{parameterised} filters were activated with the indicated parameters:
\begin{itemize}
\detokenize{$<' '.join(["\\item " + d['description'] + ' (' + next(v for v in Run1ContentVars['ConfigVars'].values() if isinstance(v, dict) and v.get('name') == 'ISMN')['parametrised_filters'][i]['parameters'].replace('_', '\\_').replace(',', ', ') + ')' for i, d in enumerate(next(v for v in Run1ContentVars['ConfigVars'].values() if isinstance(v, dict) and v.get('name') == 'ISMN')['parametrised_filters_description'])])>$}
\end{itemize}
\vspace{1em}

\noindent A list of references and acknowledgements for all used \ac{ISMN} networks is given in Table \ref{tab:ismn_refs}.

% ---------------------------------------------------------------------------------------------------------------------

\subsection{Reference -- ERA5-Land (\detokenize{$<next(v for v in Run2ContentVars['ConfigVars'].values() if isinstance(v, dict) and v.get('name') == 'ERA5-Land')['version']>$})} \label{ssec:era5land}

\noindent ERA5-Land is a global land surface reanalysis from 1950 to present, produced by the \ac{ECMWF} \citep{munozsabater2021}\footnote{available from the Copernicus Climate Data Store, \url{https://doi.org/10.24381/cds.e2161bac}}.
It was produced by replaying the land component of the \ac{ECMWF} ERA5 reanalysis \citep{hersbach2020} at an increased spatial resolution of 0.1° ($\approx$9\,km), forced by the ERA5 atmospheric fields.
Soil moisture is provided hourly and without spatial or temporal gaps for four layers; the top layer (0--7\,cm, swvl1) is used here.
In \ac{QA4SM}, the time steps at 0, 6, 12 and 18 UTC are available and regularly updated.
Only the final ERA5-Land data are used (no preliminary data), which are released with a delay of about two months.
\\[1em]
\noindent Being model-based, ERA5-Land is independent of the satellite observations, but has its own errors, e.g.\ due to errors in the precipitation forcing and the land surface model \citep{munozsabater2021}.
It is therefore used as a complementary, spatially complete reference: the comparison shows the agreement between the datasets and detects changes in the spatial coverage and quality of the products, while the accuracy of the products is assessed against the in situ data.
Observations at frozen soils are masked using the soil temperature of the same layer before the comparison.

\subsubsection{Data filters}

\noindent The following \textit{binary} filters were activated for ERA5-Land:
\begin{itemize}
\detokenize{$<' '.join([f"\\item {d['description']}" for d in next(v for v in Run2ContentVars['ConfigVars'].values() if isinstance(v, dict) and v.get('name') == 'ERA5-Land')['basic_filters_description']])>$}
\end{itemize}
